\begin{thebibliography}{9}
\phantomsection
\addcontentsline{toc}{section}{Referencias}

\bibitem{repositorio-experimentos}
Aguirre Fuentes, B. E. (2026).
\textit{EDAA-Experiments: Boletín 03, String Matching} [Código fuente].
\url{https://github.com/StackDs/EDAA-Experiments/tree/main/Boletin_03}.

\bibitem{edaa-suffix-array}
Implementación de Suffix Array, repositorio del curso EDAA.
\url{https://github.com/jfuentess/edaa/blob/main/strings/suffix_array.cpp}.

\bibitem{edaa-kmp}
Implementación de KMP, repositorio del curso EDAA.
\url{https://github.com/jfuentess/edaa/blob/main/strings/kmp.cpp}.

\bibitem{edaa-rabin-karp}
Implementación de Rabin--Karp, repositorio del curso EDAA.
\url{https://github.com/jfuentess/edaa/blob/main/strings/rk.cpp}.

\bibitem{pizzachili}
Repositorio de textos Pizza \& Chili.
\url{http://pizzachili.dcc.uchile.cl/texts.html}.

\bibitem{lovera-uhr}
Lovera, L. (LELE). (s.\,f.).
\textit{uhr: generic time performance tester} [Código fuente].
Plantilla de medición de la ayudantía, disponible en
\texttt{utils/uhr.cpp} del repositorio.
\url{https://github.com/leonardlover}.

\par\medskip
\noindent\textbf{Declaración de uso de IA generativa.}
Se utilizaron herramientas de inteligencia artificial generativa como apoyo para las implementaciones de los experimentos y para revisar la redacción técnica de este informe. Codex se utilizó como agente prioritario durante este proceso. El autor es responsable de las decisiones finales, la verificación de los resultados y el contenido presentado.

\end{thebibliography}
